\documentclass[10pt]{beamer}

\usepackage{verbatim}
%\usepackage[latin1]{inputenc}
\usepackage{amssymb,amsmath,amsfonts,url}
\usepackage{graphicx,color}
\usepackage{pgf,pgfarrows,pgfnodes,pgfautomata,pgfheaps,pgfshade}
\usepackage[utf8]{inputenc}
\usepackage[french]{babel}
\usepackage{tikz}

\usepackage[all]{xy}
\usepackage{amstext}
\usepackage{xspace}
\usepackage{mdwtab}


\mode<article> {
  \usepackage{fullpage}
  \usepackage{hyperref}
  \setjobnamebeamerversion{advancedcourse.beamer}
}





\newcommand{\alertonce}[1]{\addtocounter{beamerpauses}{-1}\alert<+>{#1}}
\newcommand{\alerttwice}[1]{\addtocounter{beamerpauses}{-1}\alert<+-+(1)>{#1}}

%\newcommand{\subsubsubsection}[1]{\noindent\textbf{#1}}

\newcommand{\cfbox}[1]{\begin{center}\fbox{#1}\end{center}}
\newcommand{\cbox}[1]{\begin{center} #1 \end{center}}

%%%%%%%%%%%%%%%%%%%%%%%%%%%%%%%%%%%%%%%%%%%%%%%%%%%%%%%%%%%%%%%%%%%%%%%%%%%%%%%%%%%%%%%%%

\graphicspath{{Graphics/}}
%\input def_0.tex

%\input{StyleSlides.tex}
\mode<presentation> {
    \usepackage{beamerthemeBoadilla}

  %\beamertemplatetransparentcovereddynamic
  \beamertemplateballitem
}



\title[IN100 Cours 9]{\textbf{Fondements informatiques I}\\
Cours 9: Algorithmes : manipuler des tableaux}

\author[]{Sorina Ionica \texttt{sorina.ionica@uvsq.fr} \\~\\Sandrine Vial  \texttt{sandrine.vial@uvsq.fr} }


\institute{}

\date{}


%\setbeamercolor{normal text}{fg=black}

\begin{document}


\frame{\titlepage}



\frame[containsverbatim]{
\frametitle{Qu'est-ce que c'est qu'un algorithme?}

\begin{alertblock}{Rappel (premier cours)} 
Un algorithme est une suite finie d'instructions et opérations permettant de résoudre un problème. 
Pour donner un algorithme: 
 
 \begin{itemize}
 \item Input: les données en entrée du problème
 \item Output : la solution du problème
 \item Les étapes qu'on doit effectuer pour résoudre
 \end{itemize}
 \end{alertblock}

\vfill

Exemple concret :  Dans un tableau de notes des étudiants de L1, on veut calculer la note la plus petite et la plus grande. \\

\vfill

Ensuite on voudrait faire afficher le tableau (la liste) des notes  à un examen,  triée dans l'ordre croissante. 


}


\frame[containsverbatim]{
\frametitle{Pour quoi?}

\textcolor{blue}{La bonne pratique:} avant de commencer à coder, on réfléchit sur papier!\\

\vspace{0.5 cm}

En langage humain, on précise les étapes de notre algorithme (pseudocode). \\

\vspace{0.5 cm}

On réfléchit également à la structure du programme :  \\

\vspace{0.3 cm}
\begin{center}
Quelles variables? De quel type? Quelles fonctions? \\
\end{center}
\vspace{0.5 cm}

Pour chaque fonction on précise les paramètres en entrée et en sortie. \\

}



\frame[containsverbatim]{
\frametitle{Minimum d'un tableau}


On veut calculer le minimum d'un tableau et renvoyer son indice. 


\begin{alertblock}{Algorithme}
\begin{verbatim} 
Input: un tableau L d'entiers
Output: l'indice du plus petit entier dans la liste

Au début on initialise m=0 (premier indice du tableau)
Pour tout i allant de 1 à longueur(L)-1
     Si L[j] < L[m] alors
            m <- j 
On renvoit m
\end{verbatim}
\end{alertblock}

}



\frame[containsverbatim]{
\frametitle{Minimum d'un tableau}


En Python cela donne: 


\begin{exampleblock}{}
\begin{verbatim} 
def minimum(L):
    m = 0
    for j in range(1, len(L)):
        if L[j] < L[m]:
            m = j
    return m
\end{verbatim}
\end{exampleblock}

\vspace{0.5 cm}

Et pour calculer le maximum? 

}


\begin{frame}[fragile]{Tri par séléction}



\begin{exampleblock}{Algorithme}
\begin{verbatim}
Input: le tableau L
Output: le tableau L trié

Pour tout i allant de 0 à longueur(L)-1: 
        On calcule l'indice m du plus petit élément dans L[i:]
        On échange L[i] et L[m]
\end{verbatim}
\end{exampleblock}

\vspace{0.5 cm}

\pause

\small{
Exécution sur le tableau  $L = [5, 1, -8, 2, -4, 7]$ :

\begin{itemize}
\item $i=0$    minimum(L)=2 alors $L = [-8, 1, 5, 2, -4, 7]$ 
\pause
\item $i=1$    minimum(L[1:])=4 alors $L = [-8, -4, 5, 2, 1, 7]$ 
\pause
\item $i=2$    minimum(L[2:])=2 alors  $ L = [-8, -4, 1, 2, 5, 7]$ 
\pause
\item $i=3$    minimum(L[3:])=3 alors  $ L = [-8, -4, 1, 2, 5, 7]$ 
\pause
\item $i=4$    minimum(L[4:])=4 alors  $ L = [-8, -4, 1, 2, 5, 7]$ 
%\item $i=5$    minimum(L[5:])=3 alors  $ L = [-8, -4, 1, 2, 5, 7]$ 
\end{itemize}
}
\end{frame}

\frame[containsverbatim]{
\frametitle{Tri par sélection en Python}


\medskip
\begin{alertblock}{}
\begin{verbatim}
def tri_selection(L):
    for i in range(len(L)):
        m = minimum(L[i:])
        L[i], L[m+i] = L[m+i], L[i] # échanger L[i] et L[m]
    return L

L = [5, 1, -8, 2, -4, 7]
print(tri_selection(L))
\end{verbatim}
\end{alertblock}


}



\frame[containsverbatim]{
\frametitle{Comment créer un tableau trié directement}

Si le tableau est crée au fur et à mesure, on insère les éléments dans le bon ordre.


\begin{alertblock}{Algorithme}
\begin{verbatim}
Input : un tableau Tab trié, un élément t
Output:  le tableau trié Tab modifiée après rajout de t


On initialise i<-0. 
taille<-longueur(Tab)
Tant que i<taille and t>l[i]
       On avance dans la liste i=i+1
On décale les éléments Tab[i], ..., Tab[taille-1] d'un cran     
On place l'élément t à l'indice i dans tab (Tab[i]=t)
On renvoie Tab       
\end{verbatim}

\end{alertblock}



}



\begin{frame}[fragile]{Comment créer un tableau trié directement}



\begin{alertblock}{Algorithme}
\begin{verbatim}
On initialise i=0. 
taille<-longueur(Tab)
Tant que i<taille and t>l[i]
      On avance dans la liste i=i+1
On décale les éléments Tab[i], ..., Tab[taille-1] d'un cran     
On place l'élément t à l'indice i dans Tab (Tab[i]=t)
On renvoie Tab   
\end{verbatim}

\end{alertblock}

\small{
Exécution pour construire  un tableau trié à partie de la séquence  5, 1, -8, 2, -4, 7:

\pause
\begin{itemize}
\item $i=0$  alors    $Tab=[5]$
\pause
\item $i=1$ alors     $Tab=[1,5]$
\pause
\item $i=2$ alors   $Tab=[-8,1,5]$
\pause
\item $i=3$ alors    $Tab=[-8,1,2,5]$
\pause
\item $i=4$ alors      $Tab=[-8,-4,1,2,5]$
\pause
\item $i=5$ alors     $Tab=[-8,-4,1,2,5,7]$ 
\end{itemize}
}

\end{frame}




\frame[containsverbatim]{
\frametitle{Insertion dans une liste triée en Python}


\begin{exampleblock}{}
\begin{verbatim}
def insertion(t,l):
    i=0
    taille=len(l)
    while (i<taille and t>l[i]):
        i+=1
    if (i<taille):    
        l[i+1:]=l[i:]
        l[i]=t
    else : 
        l.append(t)



insertion(5,l)
insertion(10,l)
print(l)
\end{verbatim}
\end{exampleblock}



}



\frame[containsverbatim]{
\frametitle{Suite de Fibonnaci}


La suite de Fibonacci est définie par la relation suivante :
\[
F_0 = 0, \quad F_1 = 1, \quad \text{et pour tout } n \ge 2, \quad F_n = F_{n-1} + F_{n-2}.
\]

%\begin{itemize}
%\item[1.] Écrire une fonction \texttt{fibonacci} qui prend en entrée un entier $n$ et qui calcule le $n$-ième nombre de Fibonnaci. 
%\item[2.] Écrire une fonction  \texttt{est\_fibonacci}  qui prend en entrée un entier $x$. La fonction devra  générer la suite jusqu’à dépasser la valeur \texttt{x}, puis vérifier si \texttt{x} y apparaît.
%\end{itemize}

On veut calculer toutes les éléments de la suite de Fibonnaci, jusqu'au $n$-ième. 

\vspace{0.5 cm}


\begin{alertblock}{Algorithme}
\begin{verbatim}
Input: n, valeurs initiales f0=0, f1=1
Output : Le tableau Fib contenant les nombres de Fibonnaci


Fib=[0,1]
i=2
Tant que  (i <=n) faire: 
      On rajoute Fib[i]=Fib[i-2]+Fib[i-1] dans le tableau
      i<-i+1
Renvoyer Fib

\end{verbatim}
\end{alertblock}




}


\frame[containsverbatim]{
\frametitle{Suite de Fibonacci}

On veut calculer le $n$-ième élément de la suite de Fibonacci. \\

On remarque qu'à chaque étape on a besoin de connaître que les deux valeurs précédentes. 

\begin{alertblock}{Algorithme}
\begin{verbatim}
Input: n, valeurs initiales a=0, b=1
Output : n-ième valeur dans la suite

i=2
Tant que  (i <= n) faire: 
      a,b<- b, a+b
      i<-i+1
Renvoyer a

\end{verbatim}
\end{alertblock}

}



\frame[containsverbatim]{
\frametitle{Suite de Fibonacci}



\begin{exampleblock}{En Python}
\begin{verbatim}
def fibonacci(n):
    if n < 0:
       print("n doit être >= 0")
    a, b = 0, 1
    for _ in range(n):
        a, b = b, a + b
    return a

\end{verbatim}
\end{exampleblock}


\vspace{0.5 cm}

\textbf{\`A retenir:} Un programme correct n'est n'est pas juste un programme qui fonctionne mais aussi: 
\begin{itemize}
\item qui utilise le moins de mémoire possible
\item qui résout le problème avec le moins d'opérations possibles!
\end{itemize}


}






\end{document}
